\chapter{Quantum Circuits of Schr\"odingerisation}\label{ch:sec06}

\section{Representation of finite-difference operators}\label{sec:sec06-fd}

Let $\Omega := [0, L]$, $N_x = 2^{n_x}$, and $h = \tfrac{L}{N_x}$. A discrete
function $u_i$ is defined on $x_i = ih$, $i = 0, 1, \dots, N_x - 1$, and encoded as
\begin{equation}
  \ket{u} := \sum_{j=0}^{2^{n_x}-1} u_j\ket{j}.
\end{equation}

Define the shift operators
\begin{align}
  S^- &:= \sum_{j=1}^{2^{n_x}-1} \ket{j-1}\bra{j} = \sum_{\bar{j}=1}^{n_x} I^{\otimes(n_x - \bar{j})} \otimes \sigma_{01} \otimes \sigma_{10}^{\otimes(\bar{j}-1)} \triangleq \sum_{\bar{j}=1}^{n_x} S^-_{\bar{j}}, \\
  S^+ &:= (S^-)\dg = \sum_{j=1}^{2^{n_x}-1} \ket{j}\bra{j-1} = \sum_{\bar{j}=1}^{n_x} I^{\otimes(n_x - \bar{j})} \otimes \sigma_{10} \otimes \sigma_{01}^{\otimes(\bar{j}-1)} \triangleq \sum_{\bar{j}=1}^{n_x} S^+_{\bar{j}},
\end{align}
which act as $S^+\ket{j} = \ket{j+1}$ ($j = 0, \dots, 2^{n_x}-2$) and
$S^-\ket{j} = \ket{j-1}$ ($j = 1, \dots, 2^{n_x}-1$), where
$\sigma_{01} := \ket{0}\bra{1}$, $\sigma_{10} := \ket{1}\bra{0}$,
$\sigma_{00} := \ket{0}\bra{0}$, and $\sigma_{11} := \ket{1}\bra{1}$. For example,
for $n_x = 1$, $S^- = \ket{0}\bra{1} = \sigma_{01} = \begin{psmallmatrix} 0 & 1 \\ 0 & 0 \end{psmallmatrix}$;
for $n_x = 2$,
\begin{equation}
  S^- = \ket{0}\bra{1} + \ket{1}\bra{2} + \ket{2}\bra{3}
  = I \otimes \sigma_{01} + \sigma_{10} \otimes \sigma_{01} \ \text{(in matrix form)};
\end{equation}
and for $n_x = 3$,
\begin{equation}
  S^- = \sum_{j=1}^{7}\ket{j-1}\bra{j} = I^{\otimes 2} \otimes \sigma_{01} + I \otimes \sigma_{10} \otimes \sigma_{01} + \sigma_{10} \otimes \sigma_{01}^{\otimes 2}.
\end{equation}

The forward and backward finite-difference (FD) operators are
\begin{equation}
  (D^+ u)_j = \frac{u_{j+1} - u_j}{h}, \qquad (D^- u)_j = \frac{u_j - u_{j-1}}{h}, \qquad j = 0, 1, \dots, N_x - 1.
\end{equation}
With periodic boundary conditions,
\begin{equation}
  D_p^+ = \frac{S^- - I^{\otimes n_x} + \sigma_{10}^{\otimes n_x}}{h}, \qquad
  D_p^- = \frac{I^{\otimes n_x} - S^+ - \sigma_{01}^{\otimes n_x}}{h},
\end{equation}
\begin{equation}
  D_p^{\pm} = \frac{S^- - S^+ - \sigma_{01}^{\otimes n_x} + \sigma_{10}^{\otimes n_x}}{2h}, \qquad
  D_p^{\Delta} = \frac{S^- + S^+ - 2I^{\otimes n_x} + \sigma_{01}^{\otimes n_x} + \sigma_{10}^{\otimes n_x}}{h^2}.
\end{equation}
Similarly, for Dirichlet boundary conditions,
\begin{equation}
  D_D^+ = \frac{S^- - I^{\otimes n_x}}{h}\ \text{(right side)}, \qquad
  D_D^- = \frac{I^{\otimes n_x} - S^+}{h}\ \text{(left side)},
\end{equation}
and
\begin{equation}
  D_D^{\Delta}\ket{u} = \frac{1}{h}(S^- - I^{\otimes n_x})\sum_{j=0}^{2^{n_x}-1} u_j\ket{j}
  = \sum_{j=0}^{2^{n_x}-1} \frac{u_{j+1} - u_j}{h}\ket{j}, \qquad u_{2^n} = 0.
\end{equation}

\begin{lemma}\label{lem:sec06-decomp}
Given an operator of the form $S = \ket{a}\bra{b} + \ket{b}\bra{a}$ with
$\ket{a}, \ket{b} \in \C^{2^n}$ and $\braket{a|b} = 0$, it can be decomposed as
\begin{equation}
  S = \frac{\ket{a}+\ket{b}}{\sqrt{2}}\frac{\bra{a}+\bra{b}}{\sqrt{2}} - \frac{\ket{a}-\ket{b}}{\sqrt{2}}\frac{\bra{a}-\bra{b}}{\sqrt{2}}
  \triangleq \ket{a'}\bra{a'} - \ket{b'}\bra{b'},
\end{equation}
where $\braket{a'|b'} = 0$. There exists a unitary matrix $B$ such that
$B\ket{0}\ket{1}^{\otimes(n-1)} = \ket{a'}$ and $B\ket{1}^{\otimes n} = \ket{b'}$,
and $S$ can be written as
\begin{equation}
  S = B\bigl(Z \otimes \ket{1}\bra{1}^{\otimes(n-1)}\bigr)B\dg.
\end{equation}
\end{lemma}

\begin{remark}
For $S = \ee^{\ii\lambda}\ket{a}\bra{b} + \ee^{-\ii\lambda}\ket{b}\bra{a}$
(with $a \to a$, $\ee^{-\ii\lambda}b \to b$), the Bell basis is
\begin{equation}
  \ket{a'} = \frac{\ket{a} + \ee^{-\ii\lambda}\ket{b}}{\sqrt{2}}, \qquad
  \ket{b'} = \frac{\ket{a} - \ee^{-\ii\lambda}\ket{b}}{\sqrt{2}},
\end{equation}
and $B = \ket{a'}\bra{0} + \ket{b'}\bra{1}$. Moreover,
\begin{equation}
  U_j(\lambda) := \Bigl(\prod_{m=1}^{j-1} \mathrm{CNOT}_m^j\Bigr) P_j(\lambda) H_j, \qquad
  P_j(\lambda) = \begin{psmallmatrix} 1 & 0 \\ 0 & \ee^{\ii\lambda} \end{psmallmatrix},
\end{equation}
where $\mathrm{CNOT}_m^j$ is the CNOT gate acting on the $m$-th qubit controlled by
the $j$-th qubit.
\end{remark}

\section{The heat equation}\label{sec:sec06-heat}

Consider the heat equation with Dirichlet boundary conditions,
\begin{equation}
  \begin{cases}
    \partial_t u(t, x) = a\Delta u(t, x), & (a > 0),\ x \in \Omega \subset \R^d, \\
    u(0, x) = f(x). &
  \end{cases}
\end{equation}
The $d$-dimensional heat equation can be solved by the central-difference method,
\begin{equation}
  \frac{\dd}{\dd t} u(t) = a D_D^{\Delta} u(t), \qquad u(0) = u_0,
  \qquad u := [u_{j_1 \cdots j_d}]_{1 \le j_i \le N_x - 1}.
\end{equation}

\paragraph{Schr\"odingerisation.}
Applying the warped phase transformation $v(t, x, p) = \ee^{-p} u(t, x)$ for
$p > 0$,
\begin{equation}
  \partial_t v(t, x, p) = a\Delta v(t, x, p) = -a\partial_p \Delta v(t, x, p),
\end{equation}
and extending to $p < 0$ with initial data $v(0, x, p) = \ee^{-p} u(0, x)$.
Applying $\mathcal{F}_p$, the Fourier transform over $p$, one obtains
\begin{equation}
  \partial_t \hat{v}(t, x, \eta) = \ii\eta\, a\Delta\hat{v}(t, x, \eta),
  \qquad \hat{v}(t, x, \eta) = \mathcal{F}_p v(t, x, p).
  \label{eq:sec06-star}
\end{equation}
Since $\eta a\Delta$ is self-adjoint under Dirichlet boundary conditions,
\eqref{eq:sec06-star} is a Schr\"odinger equation.

For the discretization, e.g.\ in one dimension, $v$ and $\hat{v}$ are discretized
as $v(t) := [v_{j,k}(t)]_{j,k}$ and $\hat{v}(t) := [\hat{v}_{j,k}(t)]_{j,k}$, with
$v_{j,k}(t) := v(t, x_j, p_k)$ and $\hat{v}_{j,k}(t) := \hat{v}(t, x_j, \eta_k)$,
and initial condition $v(0) := u(0) \otimes \vec{P}$,
$\hat{v}(0) := \mathcal{F}_p v(0) = u(0) \otimes \vec{\eta}$, where
\begin{equation}
  \vec{P} := [\ee^{-\abs{P_0}}, \ee^{-\abs{P_1}}, \dots, \ee^{-\abs{P_{N_p - 1}}}], \qquad
  \vec{\eta} := \Bigl[\tfrac{2}{\eta_0^2 + 1}, \tfrac{2}{\eta_1^2 + 1}, \dots, \tfrac{2}{\eta_{N_p-1}^2 + 1}\Bigr],
\end{equation}
with $-\pi R = P_0 < \cdots < P_{N_p} = \pi R$, $\Delta p = \tfrac{2\pi R}{N_p}$,
$N_p = 2^{n_p}$, and Fourier variable $\eta_k = (k - \tfrac{N_p}{2})/R$,
$k = 0, 1, \dots, N_p - 1$.

The Hamiltonian is discretized using the central-difference method as follows:
\begin{align}
  H_{\mathrm{heat}} &= a D_D^{\Delta} \otimes D_\eta
    = \frac{a}{h^2}\bigl(S^- + S^+ - 2I^{\otimes n_x}\bigr) \otimes \diag(\eta_0, \dots, \eta_{N_p - 1}) \notag\\
  &= \sum_{k=0}^{N_p - 1} \frac{a\eta_k}{h^2}\Bigl[\sum_{j=1}^{n_x}(S_j^- + S_j^+) - 2I^{\otimes n_x}\Bigr] \otimes \ket{k}\bra{k}
    = \sum_{k=0}^{N_p - 1}\Bigl(k - \frac{N_p}{2}\Bigr) H_0 \otimes \ket{k}\bra{k},
\end{align}
with $H_0 := \gamma_0\bigl[\sum_{j=1}^{n_x}(S_j^- + S_j^+) - 2I^{\otimes n_x}\bigr]$
and $\gamma_0 = \tfrac{a}{h^2 R}$. In the $d$-dimensional case,
\begin{equation}
  H_{\mathrm{heat}} = a\sum_{\alpha=1}^d (D_D^{\Delta})_\alpha \otimes D_\eta
  = \sum_{k=0}^{N_p-1}\Bigl(k - \frac{N_p}{2}\Bigr)\sum_{\alpha=1}^d (H_0)_\alpha \otimes \ket{k}\bra{k},
\end{equation}
where $(\cdot)_\alpha = I^{\otimes(d-\alpha)n_x} \otimes\, \cdot\, \otimes I^{\otimes(\alpha-1)n_x}$.

\section{Quantum circuit construction}\label{sec:sec06-circuit}

We construct the explicit quantum circuit for the Hamiltonian evolution operator
$U_{\mathrm{heat}}(\tau) := \exp(\ii H_{\mathrm{heat}}\tau)$:
\begin{align}
  U_{\mathrm{heat}}(\tau)
  &= \exp\Bigl(\ii\tau\sum_{k=0}^{N_p - 1}\Bigl(k - \tfrac{N_p}{2}\Bigr)\sum_{\alpha=1}^d (H_0)_\alpha \otimes \ket{k}\bra{k}\Bigr) \notag\\
  &= \sum_{k=0}^{N_p - 1}\Bigl(\exp\Bigl(\ii\tau\sum_{\alpha=1}^d (H_0)_\alpha\Bigr)\Bigr)^{k - \frac{N_p}{2}} \otimes \ket{k}\bra{k}
   = \sum_{k=0}^{N_p - 1}\Bigl(\prod_{\alpha=1}^d (U_0(\tau))_\alpha\Bigr)^{k - \frac{N_p}{2}} \otimes \ket{k}\bra{k},
\end{align}
where $U_0(\tau) := \exp(\ii H_0\tau)$ and
$\exp(\sum_k A_k \otimes \ket{k}\bra{k}) = \sum_k \exp(A_k) \otimes \ket{k}\bra{k}$,
so that $\exp(\sum_\alpha (A)_\alpha) = \prod_\alpha \exp(A_\alpha)$.

To simulate $H_0 := \gamma_0[\sum_{j=1}^{n_x}(S_j^- + S_j^+) - 2I^{\otimes n_x}]$,
consider the quantum circuit for
$\nu\sum_{j=1}^{n_x}(\ee^{\ii\lambda}S_j^- + \ee^{-\ii\lambda}S_j^+)$ with
$\lambda, \nu \in \R$. A direct calculation yields
\begin{equation}
  \ee^{\ii\lambda}S_j^- + \ee^{-\ii\lambda}S_j^+
  = I^{\otimes(n_x - j)} \otimes \bigl(\ee^{\ii\lambda}\ket{a_j}\bra{b_j} + \ee^{-\ii\lambda}\ket{b_j}\bra{a_j}\bigr),
\end{equation}
with $\ket{a_j} := \ket{0}\ket{1}^{\otimes(j-1)}$ and
$\ket{b_j} := \ket{1}\ket{0}^{\otimes(j-1)}$. Recalling \cref{lem:sec06-decomp},
one can define $B_j(\lambda)$ such that
\begin{equation}
  B_j(\lambda)\ket{0}\ket{1}^{\otimes(j-1)} = \frac{\ket{a} + \ee^{-\ii\lambda}\ket{b}}{\sqrt{2}}, \qquad
  B_j(\lambda)\ket{1}\ket{1}^{\otimes(j-1)} = \frac{\ket{a} - \ee^{-\ii\lambda}\ket{b}}{\sqrt{2}},
\end{equation}
with (Bell basis) $B_j(\lambda) := \bigl(\prod_{m=1}^{j-1}\mathrm{CNOT}_m^j\bigr)P_j(-\lambda)H_j$.
Therefore
\begin{equation}
  \ee^{\ii\lambda}S_j^- + \ee^{-\ii\lambda}S_j^+
  = I^{\otimes(n_x - j)} \otimes B_j(\lambda)\bigl(Z \otimes \ket{1}\bra{1}^{\otimes(j-1)}\bigr)B_j(\lambda)\dg.
\end{equation}
The time-evolution operator is formulated as
\begin{align}
  \exp\bigl(\ii\nu\tau(\ee^{\ii\lambda}S_j^- + \ee^{-\ii\lambda}S_j^+)\bigr)
  &= I^{\otimes(n_x - j)} \otimes B_j(\lambda)\, \mathrm{CRZ}_j^{1,\dots,j-1}(-2\nu\tau)\, B_j(\lambda)\dg \notag\\
  &\triangleq I^{\otimes(n_x - j)} \otimes W_j(\nu\tau, \lambda),
\end{align}
where
\begin{equation}
  \mathrm{CRZ}_j^{1,\dots,j-1}(\theta) := \exp(-\ii\theta Z_j/2) \otimes \ket{1}\bra{1}^{\otimes(j-1)} + I \otimes \bigl(I^{\otimes(j-1)} - \ket{1}\bra{1}^{\otimes(j-1)}\bigr)
\end{equation}
is the multi-controlled $R_Z$ gate acting on the $j$-th qubit controlled by the
$1, \dots, j-1$-th qubits (\cref{fig:sec06-wj}). Using the first-order
Lie--Trotter scheme,
\begin{equation}
  \exp\Bigl(\ii\nu\tau\sum_{j=1}^{n_x}(\ee^{\ii\lambda}S_j^- + \ee^{-\ii\lambda}S_j^+)\Bigr)
  = \prod_{j=1}^{n_x}\exp\bigl(\ii\nu\tau(\ee^{\ii\lambda}S_j^- + \ee^{-\ii\lambda}S_j^+)\bigr) + \bigo(\tau^2)
  = \prod_{j=1}^{n_x} I^{\otimes(n_x - j)} \otimes W_j(\nu\tau, \lambda).
\end{equation}

\begin{figure}[ht]
  \centering
  \begin{quantikz}
    \lstick{$q_1$}     & \qw               & \qw      & \ctrl{3}              & \qw      & \qw                & \\
    \lstick{$q_2$}     & \qw               & \qw      & \ctrl{2}              & \qw      & \qw                & \\
    \lstick{$q_{j-1}$} & \qw               & \qw      & \ctrl{1}              & \qw      & \qw                & \\
    \lstick{$q_j$}     & \gate{P(\lambda)} & \gate{H} & \gate{R_Z(-2\gamma\tau)} & \gate{H} & \gate{P(-\lambda)} &
  \end{quantikz}
  \caption{Quantum circuit for $W_j(\gamma\tau, \lambda)$.}
  \label{fig:sec06-wj}
\end{figure}

Setting $\gamma = \gamma_0$ and $\lambda = 0$ to simulate $H$, we have
$B_j(0) = \bigl(\prod_{m=1}^{j-1}\mathrm{CNOT}_m^j\bigr)H_j$. Then
$U_0(\tau) = \exp(\ii H_0\tau)$ is approximated by
\begin{align}
  U_0(\tau) &= \exp\Bigl(\ii\gamma_0\tau\sum_{j=1}^{n_x}(S_j^- + S_j^+) - 2\ii\gamma_0\tau I^{\otimes n_x}\Bigr) \notag\\
  &= \exp(-2\ii\gamma_0\tau)\exp\Bigl(\ii\gamma_0\tau\sum_{j=1}^{n_x}(S_j^- + S_j^+)\Bigr) \notag\\
  &\simeq \mathrm{Ph}(-2\gamma_0\tau)\prod_{j=1}^{n_x} I^{\otimes(n_x - j)} \otimes W_j(\gamma_0\tau, 0)
   \triangleq V_0(\tau),
\end{align}
where $\mathrm{Ph}(\theta) := \ee^{\ii\theta} I^{\otimes n_x}$ is the global phase
gate (\cref{fig:sec06-v0}).

\begin{figure}[ht]
  \centering
  \begin{quantikz}
    \lstick{$q_1$}     & \gate[1]{W_1(\gamma_0\tau,0)} & \qw & \ \ldots\ & \qw & \gate[4]{\mathrm{Ph}(-2\gamma_0\tau)} & \\
    \lstick{$q_2$}     & \qw & \gate[2]{W_2(\gamma_0\tau,0)} & \ \ldots\ & \qw & & \\
    \lstick{$\vdots$}  &     &     &           &     & & \\
    \lstick{$q_{n_x}$} & \qw & \qw & \ \ldots\ & \gate[4]{W_{n_x}(\gamma_0\tau,0)} & &
  \end{quantikz}
  \caption{Quantum circuit for $V_0(\tau)$.}
  \label{fig:sec06-v0}
\end{figure}

Ultimately, the time-evolution operator is
\begin{equation}
  U_{\mathrm{heat}}(\tau) \simeq V_{\mathrm{heat}}(\tau) := \sum_{k=0}^{N_p - 1}\Bigl(\prod_{\alpha=1}^d (V_0(\tau))_\alpha\Bigr)^{k - \frac{N_p}{2}} \otimes \ket{k}\bra{k}.
\end{equation}
Denoting $\tilde{V}_0(\tau) := \prod_{\alpha=1}^d (V_0(\tau))_\alpha$
(\cref{fig:sec06-vtilde}),
\begin{align}
  \sum_{k=0}^{N_p - 1}\tilde{V}_0^{\,k - \frac{N_p}{2}}(\tau) \otimes \ket{k}\bra{k}
  &= \tilde{V}_0^{-\frac{N_p}{2}} \otimes I^{\otimes n_p}\sum_{k=0}^{N_p-1}\tilde{V}_0^{\,(k)}(\tau) \otimes \ket{k}\bra{k}.
\end{align}
Let $k = (k_{n_p - 1}, \dots, k_0)$, so that
$\tilde{V}_0^k(\tau) = \prod_{m=0}^{n_p - 1}\tilde{V}_0^{\,k_m 2^m}(\tau)$. Then
\begin{align}
  \sum_{k=0}^{N_p - 1}\tilde{V}_0^k(\tau) \otimes \ket{k}\bra{k}
  &= \prod_{m=0}^{n_p - 1}\Bigl(\tilde{V}_0^{\,2^m}(\tau) \otimes \ket{1}\bra{1} + I^{\otimes n_x} \otimes \ket{0}\bra{0}\Bigr),
\end{align}
which is a regular matrix product in the $n_x$ registers and a tensor product in
the $n_p$ registers (\cref{fig:sec06-vheat}). The full Schr\"odingerisation
circuit is shown in \cref{fig:sec06-full}, where the measurement requires only a
projection onto $\ket{k}$, $p_k > 0$, and $\mathcal{QFT}$ ($\mathcal{IQFT}$)
denotes the (inverse) quantum Fourier transform.

\begin{figure}[ht]
  \centering
  \begin{quantikz}
    \lstick{$q^1$} & \qwbundle{n_x} & \gate{V_0(\tau)} & \qw & \ \ldots\ & \qw & \\
    \lstick{$q^2$} & \qwbundle{n_x} & \qw & \gate{V_0(\tau)} & \ \ldots\ & \qw & \\
    \lstick{$q^d$} & \qwbundle{n_x} & \qw & \qw & \ \ldots\ & \gate{V_0(\tau)} &
  \end{quantikz}
  \caption{Quantum circuit for $\tilde{V}_0(\tau)$.}
  \label{fig:sec06-vtilde}
\end{figure}

\begin{figure}[ht]
  \centering
  \begin{quantikz}
    \lstick{$q$}         & \qwbundle{n_x d} & \gate{\tilde{V}_0(\tau)} & \gate{\tilde{V}_0^{2}(\tau)} & \ \ldots\ & \gate{\tilde{V}_0^{2^{n_p-1}}(\tau)} & \\
    \lstick{$k_0$}       & \ctrl{-1}        & \qw                      & & & \\
    \lstick{$k_1$}       & \qw              & \ctrl{-1}                & & & \\
    \lstick{$k_{n_p-1}$} & \qw              & \qw                      & & \ctrl{-1} &
  \end{quantikz}
  \caption{Quantum circuit for $V_{\mathrm{heat}}(\tau)$.}
  \label{fig:sec06-vheat}
\end{figure}

\begin{figure}[ht]
  \centering
  \begin{quantikz}
    \lstick{$\ket{u(0)}$} & \qwbundle{n_x d} & \qw              & \gate[2]{V^r(\tau)} & \qw               & \qw      & \rstick{$\ket{u(r\tau)}$} \\
    \lstick{$\ket{p}$}    & \qwbundle{n_p}   & \gate{\mathcal{QFT}} &                 & \gate{\mathcal{IQFT}} & \meter{} & \rstick{$\ket{k}$}
  \end{quantikz}
  \caption{Quantum circuit for the Schr\"odingerisation method. The measurement
    requires only a projection onto $\ket{k}$ with $p_k > 0$, and $\mathcal{QFT}$
    ($\mathcal{IQFT}$) denotes the (inverse) quantum Fourier transform.}
  \label{fig:sec06-full}
\end{figure}

Given $T = r\tau$, the operator $V_{\mathrm{heat}}(\tau)$ is applied $r$ times to
give the state $\ket{\hat{v}_D(T)} := V_{\mathrm{heat}}^r(\tau)\ket{\hat{v}(0)}$.
One can perform an inverse quantum Fourier transform on it to retrieve
$\ket{v_D(T)}$. Then a measurement (the projection
$M_k := I^{\otimes n_x d} \otimes \ket{k}\bra{k}$, $p_k > 0$) is taken to select
only the $\ket{k}$ part of the state $\ket{v_D(T)}$. Since
$v(t, x, p) = \ee^{-p} u(t, x)$ for $p > 0$, one can see
\begin{equation}
  M_k\ket{v_D(T)} \approx \frac{1}{\norm{v_D(T)}}\ee^{-p_k}\mathbf{u}(T)\ket{k}
  = \frac{\ee^{-p_k}\norm{\mathbf{u}(T)}}{\norm{v(0)}}\ket{u(T)}\ket{k}.
\end{equation}
After the measurement, the output is close to $\ket{u(T)}\ket{k}$ with probability
$\ee^{-2p_k}\norm{\mathbf{u}(T)}^2 / \norm{v(0)}^2$.
